% V24 late-half mathematical and source audit; this is translator material.
\head{第二十四版专项核查：第二篇第16节后半}
\label{audit:V24-late}
\begingroup
\def\Sym{\operatorname{Sym}}
\def\gr{\operatorname{gr}}
\def\im{\operatorname{im}}
\def\Oa{\mathcal O}
\def\Pp{\mathcal P}

\audit{V24-B01}{实际边界、划去稿与名称}{已由法文原图和英译逐段核实}
本模块从法文物理第104页首句“Arrivé à ce point admirable de généralité”起，兼收英译第68页紧接可逆层图后的编者划去说明；到法文第111页原注末止。法文第104—107页有贯页划线，末线穿过推论10（iii）后结束；第107页“Notes que…”起恢复未划去正文。第112页另起第17节，英译第73页相应标题之后也不属于本轮正文。法文第103页只用于确认前后两模块的接缝。

本段包括一般性引言、增广术语、命题8、引理8.1、命题9、推论10全部条件、零点术语、奇异子预概形、常秩引理及高秩分歧建议，未用摘要替代被划去文字。法文引理8.1之前另划去“留给你来证明”一句，已经在正文译校注V24-LAT007注明，并未把它当作现行证明。

“régulier”及“quasi-régulier”分别统一为正则和拟正则；旧稿的“正规增广”改正。正文依作者第112页对零点术语的更正，采用“奇异而非超奇异的零点”；本节按极化定义的增广称“非退化二次型增广”，原“普通二次型”的名称保留于V24-LAT004。此项后续材料只作术语校勘，不据此宣称第17节已译。特征$2$中$q=X^2+YZ$的极化有$X$方向的核，但其射影零曲线光滑，足以说明两种谓词不可混同。

\audit{V24-B02}{增广二次核的类型与完备性}{原稿措辞有错，正文已作明示订正}
来源为法文104页、英译68—69页；对应正文P004—P006及V24-LAT003—005。置$R=A/I$、$V=I/I^2$。原“由$\theta^2$生成的核”把映射$\theta^2$误当生成元；全分次理想也不能直接等同于$\Sym^2(V)$中的线。结合命题8明写的$K=\ker\theta^2$，正文订正为
\[
\ker\bigl(\theta:\Sym_R(V)\to\gr_I(A)\bigr)=(K),
\qquad K=\ker\theta^2\subset\Sym_R^2(V),
\]
其中$K$为可逆子模，且由该二次张量给出的极化
$K\otimes_RV^\vee\to V$为同构。二次型增广只承接正则增广定义中的有限投射条件（a），不同时要求$\theta$为同构；法文的“a) ci-dessus”清楚，而英译的“the above”容易误读。

这一定义没有$I$-进完备、分离、$I$有限生成或$A$为Noetherian环的额外条件。下文的有限阶分次消元也不需要这些条件。相对主部$\Pp^\infty$的出现则指形式邻域的主部构造，不能反过来把完备性添入这里的一般环定义。原稿的“正则还是拟正则”的名称犹疑仍保留，不冒充作者后来的统一术语。

\audit{V24-B03}{命题8的Hilbert函数与下界问题}{两处原文数学错误已订正}
来源为法文105页、英译69页；对应P007、V24-LAT006。若$V$在一点的秩为$d\ge1$，则
\[
\operatorname{rk}\Sym^n(V)=\binom{n+d-1}{d-1}.
\]
一个每个纤维都非零的二次关系给出各次局部裂正合列
\[
0\longrightarrow K\otimes\Sym^{n-2}(V)
\longrightarrow\Sym^n(V)
\longrightarrow(\Sym(V)/(K))_n\longrightarrow0,
\]
故正确秩为$h_d(n)=\binom{n+d-1}{d-1}-\binom{n+d-3}{d-1}$，负次数的对称幂按零计。原稿的$\binom{n+d}{d}-\binom{n-2+d}{d}$在$n=1$就给出$d+1$而非$d$，并非纯粹的英文排字错误。

原稿又说“等于（或只需$\ge$）”，并由此断言$K$为可逆模；下界版本不成立。反例为$A=k[T_1,\ldots,T_d]$、$I=(T_1,\ldots,T_d)$：所有分次分量投射且秩至少为$h_d(n)$，但$K=0$。正文保留等号版本，脚注解释能保留的较弱变体。

\textbf{订正后命题的独立补证。} 若为非退化二次型增广，上述正合列和引理8.1给出各次有限投射及秩等式。反之，由第1项及$n=2$的秩等式，满射$\Sym^2(V)\to\gr_I^2(A)$的核$K$是秩一有限投射直和因子。第3项给其极化非退化，于是各纤维的局部生成二次式非零。引理8.1说明$Q_n=(\Sym(V)/(K))_n$有限投射且秩为$h_d(n)$。自然满射$Q_n\to\gr_I^n(A)$的两端有限投射且逐点同秩，故为同构，所有次数的关系均由$K$生成。

如果预先另加“$K$为可逆且非退化”，那么第2项仅要求$\operatorname{rk}\gr_I^n(A)\ge h_d(n)$也足够：满射$Q_n\to\gr_I^n(A)$同时给相反的秩不等式，继而同秩。或者只在$n=2$保留等号、其他次数用下界，也足以先保证$K$的可逆性。正文没有把这个额外前提悄悄藏在括号里。

\audit{V24-B04}{引理8.1的次数、正秩与裂单射}{公式由运算订正，难辨笔画仍明确记录}
来源为法文105页原图及3倍局部渲染、英译69页；对应P009—P010、V24-LAT008—009。法稿的两处次数上标确有叠写，放大后仍不能把最终字样当作无疑辨读。正文选择目标次数$n$，故写乘二次式的映射$\Sym^{n-2}(M)\to\Sym^n(M)$；若改以源次数记$n$，等价写成$n\to n+2$。原余核秩公式同步按命题8订正，英译把法文$n+d$另误排作$n+3$。

有限投射$M$必须在各点有正秩，才能从“极化同构”推出二次式的各纤维非零。若$B$是域、$M=0$，则零极化$0\to0$虽为同构，但$n=2$时乘零映射$B\to0$非单射。正文排除该情形；真正最小条件是对每个$x\in\Spec(B)$都有$\phi(x)\ne0$。

\textbf{独立补证。} 逐点取剩余域后，对称代数为多项式环，乘非零齐次式单射。有限投射源、目标在局部可取基，纤维上满列秩保证某个满列阶子式在邻域可逆。可逆行列变换遂把矩阵化为$\binom{I}{0}$，证明局部裂单及余核有限局部自由。这不是在一般情形中把逐纤维单射不加条件地提升为单射，而是使用有限投射矩阵的可逆子式。相同论证对任意正次数$k$的逐纤维非零齐次式成立，也保留了作者括号内的推广建议。原稿说“$B$是一个域的谱”，正文按环的类型改为“$B$是一个域”。

\audit{V24-B05}{命题9的自然二次类、正秩条件和完整补证}{补足必要范围，并消除旧稿留下的投射性条件缺口}
来源为法文105末—106页、英译69—70页；对应P011—P012、V24-LAT010—013。记$R=A/I$、$V=I/I^2$、$\bar E=E/IE$、$H=\im\bar u$、$W=V/H=J/J^2$。对$k\in K=\ker\bar u$选提升$e\in E$，则$u(e)\in I^2$；换提升只改变$u(IE)$中的元素。因此
\[
K\longrightarrow I^2/(I^3+u(IE))
\simeq\Sym_R^2(V)/(V\cdot H)
\simeq\Sym_R^2(W)
\]
定义良好。最后的识别使用对称代数对商模的性质，不预先假设$W$投射。法文先完整列完（i）、（ii）再给此解释，正文已恢复该次序；英译中的断裂商式和漏商号不再照抄。

\textbf{零维反例和范围。} 原稿未写$1\le m\le n$。取$A=k[t]_{(t)}$、$I=(t)$、$E=A^2$、$u(a,b)=at$，则$n=1,m=2$、$W=0$、$K=k$；极化$K\otimes W^\vee\to W$是同构，但$B=k$的增广二次核为零，不是可逆模。故原（ii）在零维情形不能推出（i）。正文补入$1\le m\le n$。也可以另行要求二次类逐纤维非零；本稿采用与后续几何维数一致的正秩范围。

\textbf{从非退化自动得到$W$投射。} 在$\Spec(R)$上局部平凡化$K$，令$q\in\Sym^2(W)$为相应二次张量。其极化$\beta:W^\vee\to W$是有限个求值同态与有限个向量的和，故分解为
\[
W^\vee\xrightarrow{a}R^s\xrightarrow{b}W.
\]
若$\beta$为同构，则$b(a\beta^{-1})=\operatorname{id}_W$，所以$W$是有限自由模的直和因子。随后$0\to H\to V\to W\to0$裂开，$H$有限投射；再由$0\to K\to\bar E\to H\to0$裂开，得$\operatorname{rk}H=m-1$和$\operatorname{rk}W=n-m+1$。旧稿只给出假设$W$投射后的条件性重建，本版补上这一步，未把“核是线丛”单独当作恒秩依据。例如$R=k[t]$上的映射$R^2\to R^2$、$(a,b)\mapsto(ta,0)$的核自由秩一，但在$t=0$处纤维映射为零，余核有挠。这个旧反例仍只反驳“核可逆就有纤维恒秩”的中间推理，不声称满足完整（ii）。

\textbf{有限阶消元事实。} 设$g\in I^r$且其初始项$\operatorname{in}_r(g)$为$\gr_I(A)$中的非零因子。若$ag\in I^N$、$N\ge r$，而$a\notin I^{N-r}$，则存在次数$j<N-r$使$a\in I^j\setminus I^{j+1}$；初始项的乘积在$j+r$次非零，与$ag\in I^N$矛盾。因此
\[
(g)\cap I^N=gI^{N-r}\quad(N\ge r),
\qquad
\gr(A/(g))\simeq\gr_I(A)/(\operatorname{in}_r(g)).
\]
低于$r$的次数不受影响。这里的分次商用$I$在$A/(g)$中的像定义。整个论证只考察有限多个次数，不要求滤层分离或完备；也不由此断言$g$在未分离的$A$中本身为非零因子。

\textbf{（ii）推出（i）。} 沿$\Spec(R)$局部化，并在相应的$A$局部环中提升裂开的基；前$m-1$个$u(\phi_i)$的初始项可取为$V$中独立的一次变量，余下一个初始项为零。依次用上述事实消去前$m-1$个方程，得一个环$A'$，其伴随分次代数为$\Sym_R(W)$，而余下方程$f$属于平方滤层。$f$的二次类正是$q$。由正秩和非退化，各纤维$q$非零；引理8.1给乘$q$在各次单射。再次应用消元事实，得到$\gr_J(B)\simeq\Sym_R(W)/(q)$，且二次关系线极化非退化。

\textbf{（i）推出（ii）。} （i）已给$W$有限投射、秩$n-m+1$，所以$H$及$K$局部裂开，$K$为线丛。按相同步骤消去一次变量后，$B=A'/(f)$且$f$在平方滤层。$\Sym^2(W)$到$\gr_J^2(B)$的核由$f$的二次类生成：这是因为所有$(f)$中的元素都形如$af$，其二次类仅依赖$a\bmod I$。按（i），此核是非退化的可逆关系线。由$K$到该线的自然满射为两条线丛间满射，故为同构，得到（ii）。上述全为译者补证，不是原稿“几乎平凡”之外原有的证明文字。

\audit{V24-B06}{推论10的正则性、维数与基的选择}{法文假设已恢复，数学等价已逐项核对}
来源为法文106—107页、英译70—71页；对应P013、V24-LAT014—015。法文明确$A$为正则局部环，英文漏“regular”；秩$m$与理想$\mathfrak m,\mathfrak n$等混排按法文恢复。正文补$1\le m\le n$，理由同上；零维极化的空泛反例也会使原（iii）失效。

正则局部环的伴随分次环为$\Sym_k(\mathfrak m/\mathfrak m^2)$，故命题9可用。若$\dim_k\im\bar u=m-1$，任意前$m-1$个像线性无关的基向量均给正则参数系的一部分；商$A_{m-1}$为维数$n-m+1$的正则局部环。余下方程的二次类在该商的切空间中正是命题9的$q$。极化非退化且正秩使$q\ne0$，所以余下方程非零，且在正则局部整环$A_{m-1}$中为非零因子；所得商维数$n-m$。这证明（iii）给（ii bis）及（i）；（ii bis）给（ii），而（ii）按正则参数和二次初始项直接给（iii）。反向（i）先确定命题9所需的秩：令$r=\dim_kW$，二次型增广给$\gr_{\mathfrak n}(B)\simeq k[T_1,\ldots,T_r]/(q)$且$q\ne0$，故$\dim B=\dim\gr_{\mathfrak n}(B)=r-1$；结合$\dim B=n-m$得$r=n-m+1$。现在应用命题9得到（iii）；正则序列维数条件由Cohen--Macaulay性与生成元数控制。这些解释既保留“存在一组基”与“每一组适合的基”的区别，也保留了原四项。

英文漏掉正则性不能当成无害缩写。例如$A=\mathbb Q[[x,y]]/(x^3)$、$E=A$、$u(1)=x^2+y^2$：$A$维数一但非正则；$\bar u=0$，二次类仍非退化，故英文形式的（iii）成立；$m=1$时（ii）（a）的空组商就是非正则的$A$，因而不成立。这是对英文漏项的反例，不是对法文正则局部环版本的反例。

\audit{V24-B07}{零点术语和相对平坦}{维数条件及法文“平坦”已保留}
来源为法文107—108页、英译71页；对应P015—P018、V24-LAT016。非奇异零点既要求零概形正则，也要求$m$个方程满足预期维数；奇异而非超奇异零点亦保留推论10的正则序列条件。不能仅因二次张量非退化就删除维数条件。原来的“普通”名称依B01明示更正，非超奇异仍是两类零点的并集。

法文第108页明写$V(u)$在该点对$S$“plat”，英文误作“is that”。这一平坦性可以严格补明：光滑态射的局部环映射平坦且本质有限表示；几何纤维上的正则序列经忠实平坦下降给普通纤维上的正则序列，再由\href{https://stacks.math.columbia.edu/tag/0470}{Stacks Project，Tag 0470}逐次商去，所得商仍平坦。该引用不要求基底Noetherian。基底字母的法文涂改分别在逐纤维段统一作$S$、在后面的微分构造统一作$Y$，并不把两种相对微分混在一个公式中。

\audit{V24-B08}{微分判别、判别式线丛及正秩范围}{英译错误、原稿缺号与零秩漏洞分别订正}
来源为法文108—110页、英译71—72页；对应P019—P023及V24-LAT017—022、028。$V^1$由$V^0$上的相对雅可比矩阵的$m$阶子式切出。法文外幂的源是$\bigwedge^mE$；英文漏外幂，正文已复原。$x\in V^1$仅给秩$\le m-1$；法、英括号内直接写$m-1$，正文补“至多”。法文清楚写$x\notin V''$，英文却排$x\in V''$；正文从法文及秩判据改正。所有$V^1$按法文统一，不把英译形似$V'$的字形另立概形。

二阶判别以$1\le m\le d$为范围，使$r=d-m+1>0$。原稿若容许$m=d+1$，例$X=\mathbb A^1_k$、$u=(t,0)$给$V^1=V(t)$、$V''=\varnothing$、$W=0$，空行列式为$1$，误把零点判成非超奇异，然而它不满足两个方程的预期余维。对$m>d$，本文只保留前面的$V^1$零理想构造，不使用下文二次判别，也不对全部零维例外另作分类。此前单方程定义1在零维情形有自己的字面例外，见前半模块V24-EAR034；原稿的不同定义在这条边界上并不自动兼容，因此本次只明示所证正秩范围。

在开层$U=V^1\setminus V''$上，$K$为线丛、$W$秩$r$。极化的行列式为
\[
\det\bigl(K\otimes W^\vee\to W\bigr)
\in\Gamma\bigl(U,K^{-r}\otimes(\det W)^2\bigr).
\]
法文及英文均写$K^{-n+m-1}$，但本段维数记号是$d$，应为$K^{-d+m-1}$。法文在同一说明中残留的$S$按相对微分$\Omega^1_{X/Y}$统一为$Y$。

$V^2=V(D)$只先是$U$中的闭子预概形；不加条件不能称其单独在$V^1$中闭。超奇异点集是$V''\cup V^2$，原并号缺失按补集关系恢复。这里没有为该并集凭空指定一个理想层。判别式检测极化而非所有特征下射影二次零超曲面的光滑性。

\audit{V24-B09}{常秩引理对非约化基底的要求}{原有关键假设已保留并给出局部证明}
来源为法文110页、英译72页；对应P024—P025、V24-LAT023。假设$\bigwedge^m u=0$须作为结构层同态成立。在可逆$(m-1)$阶子式的开邻域，初等行列变换给块形$\operatorname{diag}(I_{m-1},C)$。各$m$阶子式等于$C$的条目乘一个单位，故$C=0$。由此像局部为秩$m-1$直和因子；源秩$m$时，核为线丛，余核也局部自由。反向由此块形直接得两个纤维秩判据。$m=1$时$\bigwedge^0u$为恒等，$u=0$亦给相同结论。

仅凭所有剩余域上的恒秩不够：$k[\varepsilon]/(\varepsilon^2)$上乘$\varepsilon$在唯一剩余域上秩零，却有非零像$(\varepsilon)$。这是删除层同态为零条件后的反例。对本节的$V^1$，该条件恰由它的最大子式零子概形结构保证，不能把它改成仅有同一底层点集的任意概形结构。

\audit{V24-B10}{高秩分歧建议及其独立补证}{保留作者未完成状态；译者补证限定明确的分歧结构}
来源为法文110—111页、英译72—73页；对应P026—P029、V24-LAT024—026。法文列出的光滑子预概形为$V^0,V^1$；英译跨页写$V^0,V'$，并非两次$V^0$。旧附记声称英译重复$V^0$没有依据，本版撤回该归责。$d$是$X/Y$相对维数，$N$是$Y/S$相对维数；两光滑子预概形的相对维数分别假定为$N+d-m$和$N-1$。这些等式和光滑性是独立假设，非下述高度上界的自动推论。

作者说“应当验证”且“我没有详细做过”，正文原样保留其建议性质。下面是译者对源端分歧结构
$V(\operatorname{Fitt}_0\Omega^1_{V^1/Y})$所作的补证，不能归入作者原证明，也不把它写成靶端分支轨迹。

由于$V''=\varnothing$，在$V^1$每一点附近有一个可逆的$(m-1)$阶相对雅可比子式。平凡化$E$、重排方程$u_1,\ldots,u_m$，并补入$r=d-m+1$个局部函数$t_1,\ldots,t_r$，可以令$u_1,\ldots,u_{m-1},t_1,\ldots,t_r$的相对微分成基，进而给$X/Y$的相对étale坐标。取
\[
T=V(u_1,\ldots,u_{m-1}),\qquad f=u_m|_T.
\]
则$T/Y$光滑，相对维数$r$，其相对坐标为$t_i$。在这个邻域，$V^0=V(f)$；雅可比矩阵的单位块消元说明$m$阶子式的理想在$T$上正是由$\partial_i f$生成的理想。因此$V^1$的完整定义理想为
\[
(f,\partial_1f,\ldots,\partial_rf),
\]
这是理想的等式，不只是相同的底层零点。

对该商取相对微分，$\mathrm df$在$V^1$上为零，其余关系是
$\mathrm d(\partial_i f)=\sum_j(\partial_j\partial_i f)\,\mathrm dt_j$，所以有表示
\[
\Oa_{V^1}^{\oplus r}\xrightarrow{H}
\Oa_{V^1}^{\oplus r}\longrightarrow\Omega^1_{V^1/Y}\longrightarrow0,
\qquad H=(\partial_i\partial_jf).
\]
故第零Fitting理想是$(\det H)$。Fitting理想不依赖表示且与换底相容，参见\href{https://stacks.math.columbia.edu/tag/07Z6}{Stacks Project，Tag 07Z6，Lemma 15.8.2及15.8.4}。

另一方面，在同一局部平凡化中，二次初始项$q$的极化矩阵也是$H$：平方单项式$cT_i^2$贡献对角元$2c$，混合项$cT_iT_j$贡献两个对称的非对角元$c$。因此不论特征，$D$与$\det H$只差线丛平凡化所给的单位因子，所生成的零理想相同。Fitting理想与线丛截面的零理想都能粘合，得到
\[
V^2=V\bigl(\operatorname{Fitt}_0(\Omega^1_{V^1/Y})\bigr)
\]
作为$V^1$（因而在原假设下亦作为$X$）的闭子预概形相等。由于$V^1\to Y$局部有限表示，其非分歧性可由相对微分模消失判别，故底层集合也相同。此论证全程未除以$2$，也未把极化非退化改成光滑二次型。

\audit{V24-B11}{最大子式退化轨迹的余维方向}{原法文与英译同错；已改为分支余维上界}
来源为法文111页原注、英译73页提前排出的段落；对应P030、V24-LAT027。原稿在$\operatorname{rk}E=m\le d=\operatorname{rk}F$时断言秩降轨迹余维$\ge d-m+1$。取$X=\Spec k[t]$、$u:\Oa_X\to\Oa_X^2$、$1\mapsto(t,0)$，轨迹为$V(t)$、余维一，而$d-m+1=2$，即得反例。零映射给整片$X$，也说明不可把预期余维当作无条件下界。

正确结论是非空退化轨迹的每个不可约分支余维至多$d-m+1$。沿各分支泛点局部化到正则局部环，最大子式理想的高度上界给此结论。已核对一手文献：D.~Eisenbud and E.~G.~Evans, Jr., \emph{A generalized principal ideal theorem}, \emph{Nagoya Mathematical Journal} 62 (1976), 41--53，Corollary~2.3（Eagon--Macaulay--Northcott），\href{https://doi.org/10.1017/S0027763000024740}{doi:10.1017/S0027763000024740}。其公式是
$\operatorname{ht}I_k(\varphi)\le(s-k+1)(t-k+1)$；取最大子式即可。该文第二节假设局部环有Cohen--Macaulay模，正则局部环满足。高度上界不推出前段所假定的预期维数等号。

\audit{V24-B12}{旧附56—62的有效内容迁移}{按条目吸收入活动正文脚注或本报告；历史源码另存}
旧文件\texttt{text/v1\_2.tex}中的后置译校说明不再重复作为另一份正文。它的有效数学内容按以下关系保留；整份旧源码的字节存档只是历史证据，不能代替这一实质迁移。
\begin{longtable}{p{.16\textwidth}p{.75\textwidth}}
\toprule
旧附记 & 新活动位置与处理\\
\midrule
\endhead
56 & 条件（a）与（a）加（b）的区别、核的二次分量、划线边界：P001、P004—P006；LAT001、003—005；B01—B02。旧注“本批未译命题8及以后”属旧阶段状态，已由本轮真实范围取代。\\
57 & 原Hilbert秩障碍、次数重建、逐纤维非零的裂单射、仅下界不足的多项式反例：LAT006、008—009；B03—B04。新增零秩反例，并把已确证的公式订正落实正文。\\
58 & 法英段次与商式、提升无关性、条件性的分次消元补证，以及“核可逆不自动给恒秩”的警示：LAT010—013；B05。补上非退化极化自动给$W$投射的有限张量论证，明确必要正秩范围，并保留无需完备/分离的有限阶证明。\\
59 & 法文正则性、英文符号混排、非正则环反例、划去范围：LAT001、014—015；B01、B06。反例仍只针对英文漏项，未归责法文正确假设。\\
60 & 基底字母、外幂源及相对平坦性：LAT016—017、020；B07—B08。给出适用任意基底的本质有限表示平坦性判据；不再留下“通常条件适用时”的未明限定。\\
61 & 非约化基底上的常秩块矩阵、判别式线丛、开层与并集区别、特征$2$区别：LAT018—023、028；B08—B09。新增$m=d+1$的空行列式反例。\\
62 & 局部étale坐标消元、Hessian表示、Fitting理想及不除以$2$的比较：LAT026；B10。仍单列为译者补证，保留原作者待验证的叙述。\\
\bottomrule
\end{longtable}
旧附记关于法102—103图、正规性维数及可逆层扭曲的内容属于前半模块，交由前半核查报告处理，不能因本模块边界而认作本轮新增或遗漏。本文的判断只覆盖列明的原文和旧附56—62，不构成对全套EGA或全部历史译文的无错认证。
\endgroup
